% ---------------------------------------------------------------------------
% Author guideline and sample document for EG publication using LaTeX2e input
% D.Fellner, v1.13, Jul 31, 2008

\documentclass{egpubl}
\usepackage{eurovis2014}

% --- for  Annual CONFERENCE
% \ConferenceSubmission   % uncomment for Conference submission
% \ConferencePaper        % uncomment for (final) Conference Paper
\STAR                   % uncomment for STAR contribution
% \Tutorial               % uncomment for Tutorial contribution
% \ShortPresentation      % uncomment for (final) Short Conference Presentation
% \Areas                  % uncomment for Areas contribution
% \MedicalPrize           % uncomment for Medical Prize contribution
% \Education              % uncomment for Education contribution
%
% --- for  CGF Journal
% \JournalSubmission    % uncomment for submission to Computer Graphics Forum
% \JournalPaper         % uncomment for final version of Journal Paper
%
% --- for  CGF Journal: special issue
% \SpecialIssueSubmission    % uncomment for submission to Computer Graphics Forum, special issue
\SpecialIssuePaper         % uncomment for final version of Journal Paper, special issue
%
% --- for  EG Workshop Proceedings
% \WsSubmission    % uncomment for submission to EG Workshop
% \WsPaper         % uncomment for final version of EG Workshop contribution
%
 \electronicVersion % can be used both for the printed and electronic version

% !! *please* don't change anything above
% !! unless you REALLY know what you are doing
% ------------------------------------------------------------------------

% for including postscript figures
% mind: package option 'draft' will replace PS figure by a filname within a frame
\ifpdf \usepackage[pdftex]{graphicx} \pdfcompresslevel=9
\else \usepackage[dvips]{graphicx} \fi

\PrintedOrElectronic

% prepare for electronic version of your document
\usepackage{t1enc,dfadobe}

\usepackage{egweblnk}
\usepackage{cite}

% For backwards compatibility to old LaTeX type font selection.
% Uncomment if your document adheres to LaTeX2e recommendations.
% \let\rm=\rmfamily    \let\sf=\sffamily    \let\tt=\ttfamily
% \let\it=\itshape     \let\sl=\slshape     \let\sc=\scshape
% \let\bf=\bfseries

% end of prologue

% ---------------------------------------------------------------------
% EG author guidelines plus sample file for EG publication using LaTeX2e input
% D.Fellner, v1.17, Sep 23, 2010


\title[EG \LaTeX\ Author Guidelines]%
      {Human Connectome: Visualization Techniques and Methodologies}

% for anonymous conference submission please enter your SUBMISSION ID
% instead of the author's name (and leave the affiliation blank) !!
\author[G. Conte]
       {G. Conte
%        S. Spencer$^2$\thanks{Chairman Siggraph Publications Board}
        \\
% For Computer Graphics Forum: Please use the abbreviation of your first name.
         Creative Coding Lab\\
         University of Illinois at Chicago
%        $^2$ Another Department to illustrate the use in papers from authors
%             with different affiliations
       }

% ------------------------------------------------------------------------

% if the Editors-in-Chief have given you the data, you may uncomment
% the following five lines and insert it here
%
% \volume{27}   % the volume in which the issue will be published;
% \issue{1}     % the issue number of the publication
% \pStartPage{1}      % set starting page


%-------------------------------------------------------------------------
\begin{document}

% \teaser{
%  \includegraphics[width=\linewidth]{eg_new}
%  \centering
%   \caption{New EG Logo}
% \label{fig:teaser}
% }

\maketitle

\begin{abstract}
ABSTRACT

\begin{classification} % according to http://www.acm.org/class/1998/
\CCScat{Computer Graphics}{Visual Analytics}{Human Connectome}
\end{classification}

\end{abstract}





%-------------------------------------------------------------------------
\section{Introduction}

Being able to deeply understand how the brain is working is one of the main challenges in the last years among neuroscientists. With the advent and the refinement of new technologies like fMRi and diffusion MRI, doctors are able to collect and derive data about how the regions of the brain are connected. Very frequently the map of neural connections of the brain is addressed as CONNECTOME.
Visualizing these data in an effective way allows people to navigate and explore all the wired connections that are in the brain. Moreover, thanks to this kind of visualizations is possible to understand better what differences there are between healthy subjects and other people who suffer from a wide range of neuropsychiatric illnesses like bipolar, body dysmorphic disorder, Alzheimer?s disease and late-life depression. 
Many visualization tools have been proposed in the academic literature, however the vast majority of them perform 2D visualizations. However, since this research field is quite novel, there is still room for improvement.
The aim of this work is to report and survey the visualizations tools already present in the academic literature as well as to introduce which the new trends for the near future are.\\
The paper is structured as follows. In section there is a more detailed introduction about the human connectome. Then, a detailed list of common tasks follows, while in section NUMERO I describe accurately the most interesting tool. In section ... to be continued.
%-------------------------------------------------------------------------
\section{Domain}
Human Connectome has been always considered by neuroscientists a very interesting and challenging topic. However, it is only in the last few years, when more powerful and more accurate technologies took place firmly in the research area, that more detailed studies have been conducted. In fact, thanks to very advanced techniques and algorithms, it is now possible to collect data about the functional, structural and anatomical connectivity of the brain. Then, the starting points of this investigation are the data collected from fMRI and diffusion MRI. What the neuroscientists can obtain is a parcellation of the brain into smaller subregions as well as the strength of the connections, whether structural or functional, between the regions. So, the entire Connectome could be reduced to a very dense and highly connected graph. Rubinov and Sporns in \cite{complexNetworkMeasures} were the first authors that consider the human connector as a graph and in turns they described many graph-based metrics applied to the connectome.  Since as aforementioned the networks represented are highly dense, the main challenge should be addressed is that task of "creating intuitive, informative and candid images" highlighted by Margulies et al. in \cite{visualizingHumanConnectome}

\section{Main Tasks}
One of the worse mistake that can be done when designing a visual analytics tool is the willingness of addressing a spread set of tasks. So that is why is important to clearly understand and then state what the main tasks in a given research area are. This activity it is not as simple as it might seem at a first glance. So, that is why in this section I am trying to describe as much clear as possible the main goals that neuroscientist would achieve using a human connectome visualization tool.\\
Due to the high complexity of the brain network, \textit{exploration} is the the most important task. Although some of readers may argue that the exploration task is too simple, pretty obvious and too general, in my opinion, it is not so. In fact, especially when the field is quite novel and with many uncleared aspects, as it is the one we are talking about, it is extremely relevant to allow user a visualization flexibility. Flexibility in terms of level of abstraction, 


\section{Visualization Tools}
Many visualization tools have been presented in the literature. Before describing them more in details I want to report the interesting taxonomy introduced by Marguilies et al in \cite{visualizingHumanConnectome}. In fact, they identified three main categories 

\section{Future Works and new trends}


%%-------------------------------------------------------------------------
%\subsection{Language}
%
%
%
%%-------------------------------------------------------------------------
%\subsection{Margins and page numbering}
%
%
%\noindent Long captions should be set as in Figure~\ref{fig:ex1} or
%Figure~\ref{fig:ex3}.
%
%\begin{figure}[htb]
%   % an empty figure just consisting of the caption lines
%   \caption{\label{fig:ex1}
%     'Empty' figure only to serve as an example of long caption requiring 
%     more than one line. It is not typed centered but aligned on both sides.}
%\end{figure}
%
%\noindent
%Figures which need the full textwidth can be typeset as Figure~\ref{fig:ex3}.
%
%\noindent Callouts should be 9-point Times, non-boldface type. Initially
%capitalize only the first word of section titles and first-, second-, and
%third-order headings.
%
%FIRST-ORDER HEADINGS. (For example, \textbf{1. Introduction}) should be Times
%9-point boldface, initially capitalized, flush left, with one blank line
%before, and one blank line after.
%
%SECOND-ORDER HEADINGS. (For example, \textbf{2.1. Language}) should be Times
%9-point boldface, initially capitalized, flush left, with one blank line
%before, and one after. If you require a third-order heading (we discourage
%it), use 9-point Times, boldface, initially capitalized, flush left, preceded
%by one blank line, followed by a period and your text on the same line.
%
%The headline \emph{(authors / title)} must be shortened if it uses the full 
%two column width of the main text.
%There must be enough space for the page numbers. Please use ``et al.'' if 
%there are more than three authors and specify a shortened version for your title.
%
%
%
%%-------------------------------------------------------------------------
%\subsection{References}
%
%List all bibliographical references in 9-point Times, single-spaced, at the
%end of your paper in alphabetical order. When referenced in the text, enclose
%the citation index in square brackets, for example~\cite{Lous90}. Where
%appropriate, include the name(s) of editors of referenced books.
%
%For your references please use the following algorithm:
%\begin{itemize} 
%\item \textbf{one} author: first 3 chars plus year -- 
%      e.g.\ \cite{Lous90}
%\item \textbf{two}, \textbf{three} or \textbf{four} authors: first char
%      of each family name plus year --  e.g.\ \cite{Fellner-Helmberg93} 
%      or \cite{Kobbelt97-USHDR} or \cite{Lafortune97-NARF}
%\item \textbf{more than 4} authors: first char of family name from 
%      first 3 authors followed by a '*' followed by the year -- 
%      e.g.\ \cite{Buhmann:1998:DCQ} or \cite{FolDamFeiHug.etal93} 
%\end{itemize}
%
%For BibTeX users a style file \ \texttt{eg-alpha.bst} \ is available which
%uses the above algorithm.
%
%%-------------------------------------------------------------------------
%\subsection{Illustrations, graphs, and photographs}
%
%All graphics should be centered.
%
%%%%
%%%% Figure 1
%%%%
%\begin{figure}[htb]
%  \centering
%  % the following command controls the width of the embedded PS file
%  % (relative to the width of the current column)
%  \includegraphics[width=.8\linewidth]{sampleFig}
%  % replacing the above command with the one below will explicitly set
%  % the bounding box of the PS figure to the rectangle (xl,yl),(xh,yh).
%  % It will also prevent LaTeX from reading the PS file to determine
%  % the bounding box (i.e., it will speed up the compilation process)
%  % \includegraphics[width=.95\linewidth, bb=39 696 126 756]{sampleFig}
%  %
%  \parbox[t]{.9\columnwidth}{\relax
%           For all figures please keep in mind that you \textbf{must not}
%           use images with transparent background! 
%           }
%  %
%  \caption{\label{fig:firstExample}
%           Here is a sample figure.}
%\end{figure}
%
%If your paper includes images, it is very important that they are of
%sufficient resolution to be faithfully reproduced.
%
%To determine the optimum size (width and height) of an image, measure
%the image's size as it appears in your document (in millimeters), and
%then multiply those two values by 12. The resulting values are the
%optimum $x$ and $y$ resolution, in pixels, of the image. Image quality
%will suffer if these guidelines are not followed.
%
%Example 1: 
%%
%An image measures 50\,mm by 75\,mm when placed in a document. This
%image should have a resolution of no less than 600 pixels by 900
%pixels in order to be reproduced faithfully.
%
%Example 2: 
%%
%Capturing a screenshot of your entire $1024 \times 768$ pixel display
%monitor may be useful in illustrating a concept from your research. In
%order to be reproduced faithfully, that $1024 \times 768$ image should
%be no larger than 85 mm by 64 mm (approximately) when placed in your
%document.
%
%
%%-------------------------------------------------------------------------
%\subsection{Color}
%
%\textbf{Please observe:} as of 2003 publications in the proceedings of the
%Eurographics Conference can use color images throughout the paper. No
%separate color tables are necessary.
%
%However, workshop proceedings might have different agreements! 
%Figure~\ref{fig:ex3} is an example for creating color plates.
%
%%------------------------------------------------------------------------
%\subsection{Embedding of Hyperlinks / Typesetting of URLs}
%
%Due to the use of the package \texttt{hyperref} the original behavior
%of the command $\backslash$\texttt{url} from the package \texttt{url}
%is not available. To circumvent this problem we either recommend to
%use the command $\backslash$\texttt{httpAddr} from the 
%included package \texttt{egweblnk} (see below) or to replace the
%command $\backslash$\texttt{url} by the command $\backslash$\texttt{webLink} 
%-- e.g. in cases where $\backslash$\texttt{url} has been used
%widely in BibTeX-References. In the latter case we suggest to run
%BibTeX as usual and then replace all occurences of $\backslash$\texttt{url}  by
%$\backslash$\texttt{webLink}
%
%\noindent
%The provided commands for hyperlinks, in a nutshell, are:
%
%\begin{description} \itemsep 1ex
%\item [\webLinkFont $\backslash$httpAddr \{URL without leading 'http:'\}]
%      \mbox{}\\
%      e.g. \  \httpAddr{//diglib.eg.org/EG/DL/WS}
%
%\item[\webLinkFont $\backslash$ftpAddr \{URL without leading 'ftp:'\}]
%      \mbox{}\\
%      e.g. \  \ftpAddr{//www.eg.org/EG/DL/ftpupload}
%
%\item[\webLinkFont $\backslash$URL \{url\}]
%      \mbox{}\\
%      e.g. \  \URL{http://www.eg.org/EG/DL/WS}
%
%\item[\webLinkFont $\backslash$MailTo \{Email addr\}]
%      \mbox{}\\
%      e.g. \  \MailTo{publishing@eg.org}
%
%\item[\webLinkFont $\backslash$MailToNA \{emailName\}\{@emailSiteAddress\}]
%      \mbox{}\\
%      e.g. \  \MailToNA{publishing}{@eg.org}
%
%\item[\webLinkFont $\backslash$webLink\{URL without hyperlink creation\}]
%      \mbox{}\\
%      e.g. \  \webLink{http://www.eg.org/some_arbitrary_long/but_useless/URL}
%
%\end{description}
%
%
%%------------------------------------------------------------------------
%\subsection{PDF Generation}
%
%Your final paper should be delivered as a PDF document with all typefaces
%embedded. \LaTeX{} users should use \texttt{dvips} and \texttt{ps2pdf} to
%create this PDF document. Adobe Acrobat Distiller may be used in place of
%\texttt{ps2pdf}.
%
%Adobe PDFWriter is \emph{not} acceptable for use. Documents created with
%PDFWriter will be returned to the author for revision. \texttt{pdftex} and
%\texttt{pdflatex} (and its variants) can be used only if the author can
%make certain that all typefaces are embedded and images are not downsampled
%or subsampled during the PDF creation process.
%
%Users with no access to these PDF creation tools should make available a
%PostScript file and we will make a PDF document from it.
%
%
%The PDF file \emph{must not} be change protected.
%
%%------------------------------------------------------------------------
%\subsubsection*{Configuration Notes: dvips / ps2pdf / etc.}
%
%\noindent
%\texttt{dvips} should be invoked with the \texttt{-Ppdf} and \texttt{-G0}
%flags in order to use Type 1 PostScript typefaces:
%
%\begin{verbatim}
%    dvips -t a4 -Ppdf -G0 -o my.ps my.dvi
%\end{verbatim}
%
%
%\noindent
%If you are using version 7.x of GhostScript, please use the following method of invoking \texttt{ps2pdf}, in
%order to embed all typefaces and ensure that images are not downsampled or subsampled in the PDF
%creation process:
%
%\begin{verbatim}
%  ps2pdf -dMaxSubsetPct=100 \
%         -dCompatibilityLevel=1.3 \
%         -dSubsetFonts=true \
%         -dEmbedAllFonts=true \
%         -dAutoFilterColorImages=false \
%         -dAutoFilterGrayImages=false \
%         -dColorImageFilter=/FlateEncode \
%         -dGrayImageFilter=/FlateEncode \
%         -dMonoImageFilter=/FlateEncode \
%         mypaper.ps mypaper.pdf
%\end{verbatim}
%
%
%If you are using version 8.x of GhostScript, please use this method in place of the example above:
%\begin{verbatim}
%  ps2pdf -dPDFSETTINGS=/prepress \
%         -dCompatibilityLevel=1.3 \
%         -dAutoFilterColorImages=false \
%         -dAutoFilterGrayImages=false \
%         -dColorImageFilter=/FlateEncode \
%         -dGrayImageFilter=/FlateEncode \
%         -dMonoImageFilter=/FlateEncode \
%         -dDownsampleColorImages=false \
%         -dDownsampleGrayImages=false \
%         mypaper.ps mypaper.pdf
%\end{verbatim}
%
%%------------------------------------------------------------------------
%\subsubsection*{Configuration Notes: pdftex / pdflatex / etc.}
%
%\noindent
%Configuration of these tools to embed all typefaces can be accomplished by editing the \texttt{updmap.cfg} file
%to enable inclusion of the standard (or base) 14 typefaces.
%
%Linux users can run the \texttt{updmap} script to do this:
%\begin{verbatim}
%updmap --setoption pdftexDownloadBase14 true
%\end{verbatim}
%
%Windows users should edit the \texttt{updmap.cfg} files found in their TeX installation directories (one or both
%of the following may be present):
%\begin{verbatim}
%  INSTALLDIR\texmf\web2c\updmap.cfg
%  INSTALLDIR\localtexmf\miktex\config\updmap.cfg
%\end{verbatim}
%
%Ensure the value for \texttt{pdftexDownloadBase14} is "true," and then follow the instructions found here:
%\httpAddr{//docs.miktex.org/manual/} to update your MikTeX installation.
%
%%------------------------------------------------------------------------
%\subsubsection*{Configuration Notes: Acrobat Distiller}
%
%We recommend to download and install the version of the ``CMW'' Adobe Acrobat Distiller job options file
%appropriate for your operating system and version of Acrobat from the following URL:
%
%\httpAddr{//www.cadmusmediaworks.com/index2.html}\\
%in the ``(Operating System)/Applications/Distiller Settings'' folder. The ``CMW'' job options file embeds
%all typefaces and does not downsample or subsample images when creating the PDF document.
%%------------------------------------------------------------------------
%\subsection{Copyright forms}
%
%You must include your signed Eurographics copyright release form
%when you submit your finished paper. We MUST have this form before
%your paper can be published in the proceedings.
%
%%-------------------------------------------------------------------------
%\subsection{Conclusions}
%
%Please direct any questions to the production editor in charge of
%these proceedings.
%
%-------------------------------------------------------------------------

%\bibliographystyle{eg-alpha}
\bibliographystyle{eg-alpha-doi}

\bibliography{mybib}

%-------------------------------------------------------------------------
%\newpage
%
%
%\begin{figure*}[tcb]
%  \centering
%  \mbox{} \hfill
%  % the following command controls the width of the embedded PS file
%  % (relative to the width of the current column)
%  \includegraphics[width=.3\linewidth]{sampleFig}
%  % replacing the above command with the one below will explicitly set
%  % the bounding box of the PS figure to the rectangle (xl,yl),(xh,yh).
%  % It will also prevent LaTeX from reading the PS file to determine
%  % the bounding box (i.e., it will speed up the compilation process)
%  % \includegraphics[width=.3\linewidth, bb=39 696 126 756]{sampleFig}
%  \hfill
%  \includegraphics[width=.3\linewidth]{sampleFig}
%  \hfill \mbox{}
%  \caption{\label{fig:ex3}%
%           For publications with color tables (i.e., publications not offering
%           color throughout the paper) please \textbf{observe}: 
%           for the printed version -- and ONLY for the printed
%           version -- color figures have to be placed in the last page.
%           \newline
%           For the electronic version, which will be converted to PDF before
%           making it available electronically, the color images should be
%           embedded within the document. Optionally, other multimedia
%           material may be attached to the electronic version. }
%\end{figure*}

\end{document}

